% Source: 04 -Mon Oct 5 2026.pdf, page 8. Content remains in source order.
\sourcepage{04}{8}
\section{Cook's Theorem}

\textbf{THEOREM (Cook, 1971)}
\[
\mathrm{BC\text{-}SAT}\in\classNPH.
\]
The proof of Cook's theorem is complex in the details but relies on a simple intuition.


{INFORMAL PROOF}
We need to show that given an arbitrary $L\in\classNP$, $L\polyred L_{\mathrm{BC\text{-}SAT}}$.

Given that we just know $L\in\classNP$, the only thing that we know is that, $V_L(x,y)$, a polynomial time verifier for $L$, exists. We also assume to know the bounds of the verifier, k1 and k2. 

We must define a reduction function to reduce any instance of L to ann instance of BC SAT.
\[
f_{L\to\mathrm{BC\text{-}SAT}}(x):\bits^*\longrightarrow\bits^*.
\]
The reduction function builds a boolean circuit starting from $V_L$ and x. ì
\sourcefigure[0.65]{assets/04/page-008-verifier-circuit.png}
The shape of $C_{V_L,x}$ depends on $V_L$ and $x$. Its inputs ``represent'' input $y$ of $V_L$.

We will make sure that f is such that if x is in L, it will be mapped to a satisfiable circuit. Since we are creating the mapping function, we decide to only use short certificates. 

We can upper bound the number of bits that we need in association to a specific circuit. The number of bits contributed by X is fixed. We only care about certificates Y that are smaller equal than size of x to the k1. If no positive certificates of that length exist, then there can't be a positive large certificate. 

\textbf{THE REASON WHY a short Y exists is that L is is a language in NP}


